\chapter{Connectivité privée vers les plateformes crypto}\label{ch:nw:private-connectivity-to-crypto-venues}

Une plateforme qui fonctionne dans un \omterm{def:nw:trading-in-a-public-cloud:cloud}{cloud public} peut offrir à ses teneurs de marché un \omterm{def:nw:private-connectivity-to-crypto-venues:endpoint}{point de terminaison privé} à l'intérieur du réseau du fournisseur~: leurs ordres ne
passent jamais par l'internet public et, si le point de terminaison dispose d'une interface dans la même zone que leurs machines, ils ne franchissent pas non plus de frontière de zone. Les
identifiants de zone du chapitre~16 permettent à une société de vérifier si c'est le cas. La différence n'est pas négligeable~: dans le modèle de ce chapitre, le même point de terminaison servi depuis une autre zone
ajoute environ \qty{285}{\micro\second} à chaque aller-retour médian et coûte moins d'un centime de plus par million de messages.
Une connectivité privée est peu coûteuse~; une mauvaise localisation ne coûte cher qu'en temps.

Le chapitre~17 a repéré la région et la zone de la plateforme~; ce chapitre s'y connecte. Réseaux virtuels et les deux manières de les relier, l'appairage et
les \omterm{def:nw:private-connectivity-to-crypto-venues:endpoint}{points de terminaison privés}~; ce que proposent les plateformes (du \omterm{def:nw:trading-in-a-public-cloud:pg}{groupe de placement} partagé en cluster à l'\omterm{def:nw:colocation-products-and-how-they-are-sold:mmr}{interconnexion} (cross-connect) de \omterm{def:nw:colocation-products-and-how-they-are-sold:colo}{colocation} pour les rares qui fonctionnent dans des centres de données)~;
et les paliers d'accès qui ouvrent les chemins les plus rapides. La documentation des plateformes et celle du fournisseur sont citées et datées~; les latences sont une
simulation étiquetée, construite sur la mesure publiée du chapitre~16.

\section{Réseaux virtuels, appairage et points de terminaison privés}

\begin{definition}[Cloud privé virtuel, appairage de réseaux virtuels]\label{def:nw:private-connectivity-to-crypto-venues:peering}
Un \emph{cloud privé virtuel}\index{cloud privé virtuel} (VPC) est un réseau privé d'un client à l'intérieur d'un \omterm{def:nw:trading-in-a-public-cloud:cloud}{cloud public}~: son propre espace d'adressage,
ses sous-réseaux dans les zones choisies, ses règles de routage et d'accès. L'\emph{appairage de réseaux virtuels}\index{appairage de réseaux virtuels} relie deux de ces réseaux, appartenant au même client ou à des clients différents,
de sorte que leurs machines échangent du trafic par adresses privées comme s'il s'agissait d'un seul réseau.
\end{definition}

\begin{definition}[Point de terminaison privé, répartiteur de charge réseau]\label{def:nw:private-connectivity-to-crypto-venues:endpoint}
Un \emph{point de terminaison privé}\index{point de terminaison privé} est une interface réseau que le fournisseur place dans le réseau virtuel d'un client, dans les zones choisies,
par laquelle le client atteint le service d'un tiers de façon privée, sans routage entre les deux réseaux. Un \emph{répartiteur de
charge réseau}\index{répartiteur de charge réseau} est le service du fournisseur qui reçoit les connexions sur une adresse et les répartit sur des machines cibles,
au niveau de TCP ou d'UDP~; un service de point de terminaison privé est généralement exposé à travers un tel répartiteur.
\end{definition}

L'appairage et les points de terminaison résolvent des problèmes différents. L'appairage relie deux réseaux et leur permet de router l'un vers l'autre~: rapide et simple, mais les deux côtés doivent
s'accorder sur les adresses et se faire mutuellement confiance pour leurs routes. Un point de terminaison publie un service à de nombreux clients sans relier leurs réseaux~: chaque
client dispose d'une interface dans son propre réseau, et le fournisseur achemine ses paquets, par un répartiteur de charge, jusqu'au service. AWS documente
les deux~: l'appairage permet aux instances de «~communiquer entre elles comme si elles se trouvaient dans le même réseau~»~; PrivateLink relie un réseau à un
service «~comme s'ils étaient dans votre VPC~».

\begin{omfigure}
\begin{tikzpicture}[font=\footnotesize,b/.style={draw,rounded corners=2pt,minimum height=0.55cm,align=center,inner sep=3pt},>=Stealth]
\node[b,fill=omDef!15,minimum width=2.4cm] (f) at (0,0) {instance de la société\\ (zone az4)};
\node[b,fill=omThm!15,minimum width=2.6cm] (v) at (10.4,0) {passerelles de la plateforme\\ (son VPC)};
\node[b,fill=gray!10] (e) at (3.4,-1.5) {interface du point\\ de terminaison privé};
\node[b,fill=gray!10] (l) at (7,-1.5) {répartiteur de\\ charge réseau};
\draw[->,thick] (f) -- node[above,font=\scriptsize]{appairage~: routage entre les deux VPC} (v);
\draw[->,thick,omMeth] (f) |- (e); \draw[->,thick,omMeth] (e) -- (l); \draw[->,thick,omMeth] (l) -| (v);
\node[font=\scriptsize,omMeth] at (5.2,-2.25) {point de terminaison privé~: sauts supplémentaires, pas de routage entre les réseaux};
\node[b,fill=omMeth!10] (x) at (5.2,1.6) {réseau de périphérie (public)};
\draw[->,thick,gray,dashed] (f) |- (x); \draw[->,thick,gray,dashed] (x) -| (v);
\end{tikzpicture}

\omcaption{Trois chemins depuis l'instance d'une société vers une plateforme dans la même région cloud, de façon schématique~: l'appairage entre les deux réseaux virtuels, un
\omterm{def:nw:private-connectivity-to-crypto-venues:endpoint}{point de terminaison privé} (une interface dans le réseau de la société, un répartiteur de charge devant les passerelles de la plateforme) et le point de terminaison public passant par un \omterm{def:nw:where-the-crypto-venues-live:as}{réseau de périphérie}.}\label{fig:nw:private-connectivity-to-crypto-venues:paths}
\end{omfigure}

La zone du point de terminaison est le piège. Un point de terminaison dispose d'interfaces dans les zones où il est créé~; le répartiteur de charge
placé derrière lui envoie chaque connexion vers des cibles de sa zone, sauf si l'équilibrage inter-zones est activé, option que AWS documente comme désactivée par défaut
pour les \omterm{def:nw:private-connectivity-to-crypto-venues:endpoint}{répartiteurs de charge réseau}. Si les passerelles de la plateforme se trouvent dans une zone physique et l'instance de la société dans une autre, chaque paquet franchit la
frontière quelque part~; seule la question du lieu reste ouverte. Les identifiants du chapitre~16 permettent à la société de comparer la zone de son instance avec celle des
interfaces du point de terminaison avant d'envoyer le moindre ordre.

\begin{proposition}[Où se situe la frontière de zone]\label{prop:nw:private-connectivity-to-crypto-venues:zone}
Si l'instance de la société et les passerelles de la plateforme se trouvent dans des zones physiques différentes, chaque chemin entre elles franchit une frontière de zone, quel que soit le
placement du point de terminaison~; si elles se trouvent dans la même zone, un chemin ne la franchit que si le point de terminaison ou le répartiteur de charge achemine le trafic par
une autre zone. La société a donc besoin des deux identifiants (celui de son instance et celui des passerelles de la plateforme) pour savoir si un chemin dans la même zone existe.
\end{proposition}

\begin{proof}
Un chemin relie deux points~; s'ils se situent dans des zones différentes, un lien du chemin relie les zones. S'ils se situent dans une même zone, seul un saut placé dans
une autre zone (une interface de point de terminaison ou un nœud de répartiteur de charge) fait sortir le chemin de cette zone puis l'y ramène.
\end{proof}

\begin{omfigure}
\begin{tikzpicture}[font=\footnotesize,b/.style={draw,rounded corners=2pt,minimum height=0.5cm,align=center,inner sep=2pt},>=Stealth]
\foreach \x/\z in {0/az1,4.2/az2,8.4/az4}{\draw[gray,rounded corners=3pt] (\x,-1.3) rectangle (\x+3.6,1.5); \node[gray,font=\small] at (\x+1.8,1.25) {\z};}
\node[b,fill=omThm!20] (g) at (1.8,0.4) {passerelles de\\ la plateforme};
\node[b,fill=gray!10] (e1) at (1.8,-0.6) {interface du point\\ de terminaison};
\node[b,fill=gray!10] (e2) at (6.0,-0.6) {interface du point\\ de terminaison};
\node[b,fill=omDef!15] (i) at (10.2,0.4) {instance de la société};
\draw[->,thick,omMeth] (i) .. controls (8,-0.2) and (7.5,-0.6) .. (e2);
\draw[->,thick,omMeth] (e2.west) .. controls (4.0,-0.6) and (3.6,0.4) .. (g.east);
\node[font=\scriptsize,omMeth,align=center] at (6.0,0.6) {chaque paquet franchit\\ deux frontières de zone};
\end{tikzpicture}

\omcaption{Le point de terminaison dans la mauvaise zone, de façon schématique~: les passerelles de la plateforme en az1, les interfaces de son point de terminaison en az1 et en az2, l'instance de la société en
az4~; les paquets de la société entrent par l'interface en az2 et rejoignent les passerelles à travers le répartiteur de charge. Aucun placement du point de terminaison ne maintient ce chemin dans une seule zone (la \cref{prop:nw:private-connectivity-to-crypto-venues:zone})~; déplacer l'instance en az1
le permet. Les étiquettes de zone sont des identifiants, non des noms.}\label{fig:nw:private-connectivity-to-crypto-venues:zones}
\end{omfigure}

\omcode{code/firm/privlink/firm_privlink.py}{59}{77}{Aller-retour d'un chemin (le corps du modèle du chapitre~16 augmenté des sauts supplémentaires documentés), contrôle de la
zone et coûts.}\label{lst:nw:private-connectivity-to-crypto-venues:path}

\begin{dated}{2026-09}{Tarifs des points de terminaison privés}\label{dat:nw:private-connectivity-to-crypto-venues:prices}
La liste de prix d'AWS du 17 septembre 2026 facture 0.014~USD de l'heure pour chaque interface de point de terminaison dans chaque zone d'Asie-Pacifique (Tokyo), 0.01 dans l'Est des États-Unis
(N. Virginia), et 0.01~USD par gigaoctet de données traitées pour le premier pétaoctet par mois dans une région (0.006 pour les quatre suivants, 0.004 au-delà)~; les données
qui franchissent les zones sont facturées 0.01~USD par gigaoctet dans chaque sens (chapitre~16).
\end{dated}

\section{Colocation proposée par les plateformes}

\begin{dated}{2026-09}{Ce que proposent les plateformes, d'après leur documentation}\label{dat:nw:private-connectivity-to-crypto-venues:offers}
\begin{itemize}
\item Coinbase Derivatives~: production dans Equinix CH4 (350 E Cermak, Chicago), reprise après sinistre et intégration dans Equinix NY5 (Secaucus)~; accès par
  \omterm{def:nw:colocation-products-and-how-they-are-sold:mmr}{interconnexion} (cross-connect) de \omterm{def:nw:colocation-products-and-how-they-are-sold:colo}{colocation} (recommandée pour la latence la plus basse et la plus stable), par AWS PrivateLink («~certains avantages de la connectivité privée
  sans avoir à passer par une \omterm{def:nw:colocation-products-and-how-they-are-sold:mmr}{interconnexion}~»), ou par l'internet avec TLS.
\item One Trading, avec AWS (décembre 2023)~: quatre paliers d'accès testés pour les teneurs de marché, depuis l'appairage VPC à l'intérieur d'un \omterm{def:nw:trading-in-a-public-cloud:pg}{groupe de placement} en cluster partagé
  entre les comptes de la bourse et du teneur de marché (latence la plus basse), en passant par l'appairage simple et PrivateLink derrière un \omterm{def:nw:private-connectivity-to-crypto-venues:endpoint}{répartiteur de charge réseau}
  (intermédiaire, avec des sauts supplémentaires), jusqu'à l'internet à travers un \omterm{def:nw:where-the-crypto-venues-live:as}{réseau de périphérie}~; l'objectif annoncé par la bourse, des allers-retours inférieurs à \qty{200}{\micro\second}.
\end{itemize}
\end{dated}

Un \omterm{def:nw:trading-in-a-public-cloud:pg}{groupe de placement} en cluster partagé est la version cloud de la \omterm{def:nw:colocation-products-and-how-they-are-sold:colo}{colocation}~: la plateforme crée le groupe dans son compte et le partage, de sorte que les
instances du teneur de marché se retrouvent dans le même segment à forte bissection du réseau que les passerelles et le moteur d'appariement (matching engine) de la plateforme. C'est la seule
offre de la liste qui rapproche physiquement la machine du client~; les autres choisissent un chemin vers l'endroit où se trouve la machine.

\section{Interconnexions dans les plateformes hébergées en centre de données}

Certaines plateformes crypto ne fonctionnent pas du tout dans un cloud. La bourse de dérivés de Coinbase fonctionne dans un centre de données à Chicago, avec son plan de reprise après sinistre à New
Jersey, dans les mêmes bâtiments que les plateformes du chapitre~10, et vend les connexions du chapitre~9~: des \omterm{def:nw:colocation-products-and-how-they-are-sold:mmr}{interconnexions} vers le réseau de la bourse dans le bâtiment.
Sa documentation est explicite sur la hiérarchie~: des \omterm{def:nw:colocation-products-and-how-they-are-sold:mmr}{interconnexions} pour la latence la plus basse, PrivateLink pour une connectivité privée sans
\omterm{def:nw:colocation-products-and-how-they-are-sold:mmr}{interconnexion}, l'internet pour le reste. Pour une telle plateforme, le cloud est un réseau d'accès, et non le lieu d'hébergement, et le choix de la firme est celui du chapitre~9.

\section{Passerelles dédiées et passerelles des teneurs de marché, et les paliers qui les débloquent}

Les chemins les plus rapides ne sont pas ouverts à tous. Les plateformes réservent les \omterm{def:nw:trading-in-a-public-cloud:pg}{groupes de placement} partagés, les \omterm{def:nw:private-connectivity-to-crypto-venues:endpoint}{points de terminaison privés} et les passerelles dédiées aux clients au-dessus d'un
volume ou d'une commission, comme le décrit l'encadré daté du livre~3 sur l'endpoint dédié et l'alignement des paliers, et peuvent leur appliquer des limites de débit distinctes.
Pour la firme qui remplit les conditions, la question est celle de ce chapitre~: quel chemin, depuis quelle zone. Pour celle qui ne les remplit pas, elle relève du chapitre~17~:
la meilleure instance qu'elle peut trouver sur le chemin public.

\begin{table}[ht]
\centering\small
\begin{tabular}{@{}>{\raggedright\arraybackslash}p{5.0cm}>{\centering\arraybackslash}p{2.0cm}>{\centering\arraybackslash}p{2.0cm}>{\centering\arraybackslash}p{2.0cm}>{\centering\arraybackslash}p{2.0cm}@{}}
\toprule
Chemin & p50 & p99 & Coût par million & Fixe par zone \\
 & (\si{\micro\second}) & (\si{\micro\second}) & de messages (USD) & (USD par mois) \\
\midrule
Appairage, \omterm{def:nw:trading-in-a-public-cloud:pg}{groupe de placement} en cluster partagé & 150 & 220 & 0 & 0 \\
Appairage, même zone & 270 & 396 & 0 & 0 \\
\omterm{def:nw:private-connectivity-to-crypto-venues:endpoint}{Point de terminaison privé}, même zone & 310 & 436 & 0.004 & 10.22 \\
\omterm{def:nw:private-connectivity-to-crypto-venues:endpoint}{Point de terminaison privé}, autre zone & 595 & 797 & 0.012 & 10.22 \\
Point de terminaison public via un \omterm{def:nw:where-the-crypto-venues-live:as}{réseau de périphérie} & 2\,270 & 2\,396 & 0 & 0 \\
\bottomrule
\end{tabular}
\caption{Simulation des chemins documentés dans une région~: les allers-retours ajustés du chapitre~16 (le cluster partagé retenu à la valeur «~sub-150~µs~» du blog d'AWS),
plus une hypothèse de \qty{40}{\micro\second} pour les sauts du point de terminaison et du répartiteur de charge, et \qty{2}{\milli\second} pour le \omterm{def:nw:where-the-crypto-venues-live:as}{réseau de périphérie} public~; coûts issus des prix datés,
prix, pour 400~octets par message (un ordre et son accusé de réception). Données~: \texttt{nw\_private.path\_rows()}.}\label{tab:nw:private-connectivity-to-crypto-venues:paths}
\end{table}

\begin{omfigure}
\begin{tikzpicture}
\begin{axis}[ybar,width=11.5cm,height=6cm,bar width=9pt,ymode=log,ymin=100,ymax=5000,log origin=infty,xtick=data,
  symbolic x coords={peering + shared cluster,peering,endpoint same zone,endpoint other zone,public via edge},xticklabels={peering + cluster partagé,peering,terminaison même zone,terminaison autre zone,public via périphérie},
  x tick label style={font=\scriptsize,text width=1.8cm,align=center},ylabel={aller-retour (\si{\micro\second}, log)},
  label style={font=\small},grid=major,grid style={gray!25},legend style={font=\footnotesize},legend pos=north west,
  table/col sep=comma,enlarge x limits=0.12]
\addplot[fill=omDef!55,draw=omDef] table[x=path,y=p50]{figdata/networks/18-private-connectivity-to-crypto-venues/paths.csv};
\addplot[fill=omThm!45,draw=omThm] table[x=path,y=p99]{figdata/networks/18-private-connectivity-to-crypto-venues/paths.csv};
\legend{p50,p99}
\end{axis}
\end{tikzpicture}

\omcaption{Simulation~: allers-retours médians et au 99\up{e}~centile des cinq chemins du \cref{tab:nw:private-connectivity-to-crypto-venues:paths}. Le point de terminaison
dans la mauvaise zone coûte autant que le passage de l'appairage au point de terminaison, et même plusieurs fois plus~; le \omterm{def:nw:where-the-crypto-venues-live:as}{réseau de périphérie} public coûte des millisecondes. Données~:
\texttt{fig\_private.py}.}\label{fig:nw:private-connectivity-to-crypto-venues:rtt}
\end{omfigure}

\begin{omfigure}
\begin{tikzpicture}
\begin{axis}[width=11cm,height=5.4cm,xmin=0,xmax=10000,ymin=0,ymax=140,xlabel={messages par mois (millions)},x tick label style={/pgf/number format/1000 sep={\,}},
  ylabel={coût du point de terminaison (USD par mois)},tick label style={font=\footnotesize},label style={font=\small},grid=major,grid style={gray!25},
  legend pos=north west,legend style={font=\footnotesize},table/col sep=comma,scaled x ticks=false]
\addplot[omDef,thick,no markers] table[x=millions,y=same]{figdata/networks/18-private-connectivity-to-crypto-venues/monthly.csv};
\addplot[omThm,thick,dashed,no markers] table[x=millions,y=other]{figdata/networks/18-private-connectivity-to-crypto-venues/monthly.csv};
\legend{point de terminaison dans la même zone,point de terminaison dans une autre zone}
\end{axis}
\end{tikzpicture}

\omcaption{Le coût mensuel d'un \omterm{def:nw:private-connectivity-to-crypto-venues:endpoint}{point de terminaison privé} dans une zone, en fonction du volume de messages, d'après les prix datés~: 10.22~USD d'heures de point de terminaison plus le traitement des
données, et le transfert entre zones lorsque le point de terminaison dessert une autre zone. Dix milliards de messages par mois coûtent 50 ou 130~USD. Données~:
\texttt{fig\_private.py}.}\label{fig:nw:private-connectivity-to-crypto-venues:cost}
\end{omfigure}

\omcode{code/networks/18-private-connectivity-to-crypto-venues/python/nw_private.py}{32}{36}{Le résultat nommé~: ce que l'autre zone ajoute au
aller-retour et au coût par million de messages.}\label{lst:nw:private-connectivity-to-crypto-venues:wrong}

Les coûts sont secondaires~: un point de terminaison dans une zone pendant un mois coûte 10.22~USD d'heures de point de terminaison, et même dix milliards de messages ajoutent 40~USD de
traitement des données, 120 entre zones. Ce qui compte, c'est la zone. Le point de terminaison dans la mauvaise zone ajoute \qty{285}{\micro\second} à l'aller-retour médian et
\qty{360}{\micro\second} au 99\up{e}~centile, plus que toute la différence entre l'appairage et un point de terminaison dans la même zone.

\begin{method}[Se connecter en privé à une plateforme hébergée dans le cloud]\label{met:nw:private-connectivity-to-crypto-venues:connect}
\begin{enumerate}
\item Lire la documentation de connectivité de la plateforme~: les chemins qu'elle propose (\omterm{def:nw:trading-in-a-public-cloud:pg}{groupe de placement} partagé, appairage, points de terminaison, \omterm{def:nw:colocation-products-and-how-they-are-sold:mmr}{interconnexions}), à quels clients, dans quelles
  zones, à quelles limites de débit~; consigner chacun avec sa date.
\item Retrouver les identifiants de zone de vos instances et des interfaces de point de terminaison de la plateforme~; placer les instances dans la zone des passerelles de la plateforme.
\item Privilégier le chemin comportant le moins de sauts que la plateforme autorise~; vérifier que le point de terminaison dispose d'une interface dans la zone des instances et que son répartiteur de charge garde
  le trafic dans la zone.
\item Mesurer l'aller-retour de chaque chemin avec des horodatages matériels~; conserver le chemin public comme solution de repli, et connaître sa latence.
\item Chiffrer les heures de point de terminaison et le traitement des données, et les confronter à la latence~; elles seront minimes.
\end{enumerate}
\end{method}

\section{Tutoriel~: quel chemin, quelle zone}\label{tut:nw:private-connectivity-to-crypto-venues:tut}

\begin{tutorial}
\textbf{Objectif.} Comparer les chemins documentés vers une plateforme hébergée dans le cloud selon l'aller-retour et le coût, et vérifier l'alignement des zones.
\textbf{État final~:} les \cref{fig:nw:private-connectivity-to-crypto-venues:rtt,fig:nw:private-connectivity-to-crypto-venues:cost} et
le \cref{tab:nw:private-connectivity-to-crypto-venues:paths}.
\begin{tutsteps}
\item \textbf{Chemins.} \texttt{firm\_privlink.load\_paths} lit \texttt{data/paths.csv}~: chaque chemin documenté avec ses sauts supplémentaires, ses prix et ses
  sources.
\item \textbf{Allers-retours.} \texttt{rtt\_percentiles} (voir le \cref{lst:nw:private-connectivity-to-crypto-venues:path}) ajoute les sauts aux modèles
  de corps du chapitre~16.
\item \textbf{Alignement.} \texttt{aligned} compare l'identifiant de zone de l'instance à celui du point de terminaison~; \texttt{nw\_private.ALIGN} est un exemple en échec.
\item \textbf{Coûts.} \texttt{cost\_per\_million} et \texttt{monthly\_fixed}~; \texttt{monthly\_curve} trace le coût en fonction du volume.
\end{tutsteps}
\textbf{Pour aller plus loin.} Remplacer l'hypothèse de \qty{40}{\micro\second} pour les sauts du point de terminaison par une mesure~; activer l'équilibrage de charge entre zones dans le
  modèle et router une fraction des connexions vers l'autre zone.
\end{tutorial}

\section{Construction~: le planificateur de chemins privés}\label{bld:nw:private-connectivity-to-crypto-venues:privlink}

\begin{build}
\bfield{Rôle} Les chemins des machines cloud de la firme vers chaque plateforme, leurs allers-retours et leurs coûts, et s'ils restent dans une seule zone~: l'entrée des
  lignes cloud du plan du chapitre~29.
\bfield{Interface} \texttt{firm\_privlink}~: \texttt{Path}, \texttt{load\_paths}, \texttt{BODIES}, \texttt{rtt\_percentiles}, \texttt{aligned},
\texttt{cost\_per\_million}, \texttt{monthly\_fixed}~; sur \texttt{firm.cloudplan}.
\bfield{Règles} Un chemin n'existe que si la documentation d'une plateforme le propose, avec sa source et sa date~; les sauts supplémentaires sont des hypothèses énoncées tant qu'ils ne sont pas mesurés~; un
point de terminaison non aligné est signalé, jamais accepté en silence.
\bfield{Tests d'acceptation} \texttt{code/firm/privlink/tests/}~: les sources du catalogue, le contrôle d'alignement, l'ordre des médianes des chemins, et les
coûts à la main.
\bfield{Pour aller plus loin} Lire les interfaces des points de terminaison et leurs zones à partir de l'interface du fournisseur~; ajouter les limites de débit de la plateforme par chemin.
\end{build}

\begin{omsources}
\item Documentation AWS~: appairage de VPC, PrivateLink et sa tarification, répartiteurs de charge réseau~; liste de prix publique d'AWS (septembre 2026).
\item Coinbase Developer Documentation, Derivatives connectivity~; blog AWS for Industries, «~One Trading and AWS~: Cloud-native colocation for crypto
trading~» (décembre 2023).
\end{omsources}

\section{Exercices}

\begin{exercise}[$\star$]\label{exo:nw:private-connectivity-to-crypto-venues:1}
Combien coûte un \omterm{def:nw:private-connectivity-to-crypto-venues:endpoint}{point de terminaison privé} par mois dans une zone à Tokyo, avant tout échange de données~?
\end{exercise}

\begin{exercise}[$\star$]\label{exo:nw:private-connectivity-to-crypto-venues:2}
Quelle est la différence entre l'appairage et un \omterm{def:nw:private-connectivity-to-crypto-venues:endpoint}{point de terminaison privé}, du point de vue de la plateforme~?
\end{exercise}

\begin{exercise}[$\star$]\label{exo:nw:private-connectivity-to-crypto-venues:3}
Pourquoi un \omterm{def:nw:trading-in-a-public-cloud:pg}{groupe de placement} en cluster partagé rapproche-t-il le teneur de marché, alors qu'un point de terminaison ne le fait pas~?
\end{exercise}

\begin{exercise}[$\star\star$]\label{exo:nw:private-connectivity-to-crypto-venues:4}
Votre instance se trouve dans apne1-az4~; le point de terminaison de la plateforme a des interfaces dans apne1-az1 et apne1-az2. Que faites-vous~?
\end{exercise}

\begin{exercise}[$\star\star$]\label{exo:nw:private-connectivity-to-crypto-venues:5}
Pourquoi le chemin public passant par un \omterm{def:nw:where-the-crypto-venues-live:as}{réseau de périphérie} est-il plus lent que chacun des chemins privés, même à l'intérieur d'une même région~?
\end{exercise}

\begin{exercise}[$\star\star$]\label{exo:nw:private-connectivity-to-crypto-venues:6}
Une plateforme fonctionne dans un centre de données à Chicago et propose PrivateLink. Pourquoi une firme située dans le même centre de données achèterait-elle quand même une \omterm{def:nw:colocation-products-and-how-they-are-sold:mmr}{interconnexion} (cross-connect)~?
\end{exercise}

\begin{exercise}[$\star\star\star$]\label{exo:nw:private-connectivity-to-crypto-venues:7}
\emph{Programmation.} Avec \texttt{firm.privlink}, trouver le temps de saut du point de terminaison auquel la médiane d'un point de terminaison dans la même zone égale le 99\up{e}~centile de l'appairage simple.
\end{exercise}

\begin{exercise}[$\star\star\star$]\label{exo:nw:private-connectivity-to-crypto-venues:8}
\emph{Trouver la faille.} «~Nous utilisons le \omterm{def:nw:private-connectivity-to-crypto-venues:endpoint}{point de terminaison privé} de la plateforme, donc nos ordres empruntent le chemin le plus rapide disponible.~»
\end{exercise}

\section{Problème~: le point de terminaison dans la mauvaise zone}

\begin{problem}[{Problème du week-end --- ce que coûte la mauvaise zone}]\label{pb:nw:private-connectivity-to-crypto-venues:1}
Un teneur de marché atteint une plateforme par le \omterm{def:nw:private-connectivity-to-crypto-venues:endpoint}{point de terminaison privé} de celle-ci. Ses instances se trouvent dans une seule zone~; le point de terminaison la dessert depuis une autre. Les messages font
400 octets (un ordre et son accusé de réception)~; les sauts du point de terminaison du modèle prennent \qty{40}{\micro\second}.

\medskip\noindent\textbf{Partie I --- L'aller-retour.}
\begin{enumerate}
\item Quels sont les allers-retours médian et au 99\up{e}~centile par un point de terminaison dans la même zone~?
\item Et par le point de terminaison situé dans une autre zone~?
\item De combien la mauvaise zone augmente-t-elle chacun de ces temps~?
\item Comment cela se compare-t-il au passage de l'appairage à un point de terminaison dans la même zone~?
\end{enumerate}

\medskip\noindent\textbf{Partie II --- Le coût.}
\begin{enumerate}[resume]
\item Que coûtent un million de messages par chaque point de terminaison~?
\item Quelle part de l'écart provient du transfert inter-zones~?
\item Que coûtent dix milliards de messages par mois par chacun~?
\item Que coûtent les heures de point de terminaison dans une zone et dans trois~?
\end{enumerate}

\medskip\noindent\textbf{Partie III --- La correction.}
\begin{enumerate}[resume]
\item Comment la firme découvre-t-elle qu'elle se trouve dans la mauvaise zone~?
\item Quelles sont ses options si le point de terminaison de la plateforme n'a pas d'interface dans sa zone~?
\item Que lui apporterait un \omterm{def:nw:trading-in-a-public-cloud:pg}{groupe de placement} en cluster partagé~?
\item Que dit la \cref{prop:nw:private-connectivity-to-crypto-venues:zone} du déplacement du point de terminaison plutôt que de celui des instances~?
\end{enumerate}

\medskip\noindent\textbf{Partie IV --- Le verdict.}
\begin{enumerate}[resume]
\item Énoncer le \emph{résultat nommé}~: la latence ajoutée et le coût de traitement des données par million de messages d'un \omterm{def:nw:private-connectivity-to-crypto-venues:endpoint}{point de terminaison privé} desservi depuis une autre zone,
  comparés à ceux d'une même zone.
\item Quelles données relèvent de la documentation, lesquelles de la mesure, lesquelles de l'hypothèse~?
\item Pourquoi le prix de la connectivité privée n'est-il pas une raison d'hésiter~?
\item Que devrait publier la plateforme pour prévenir l'erreur~?
\item À quelle fréquence la firme doit-elle vérifier à nouveau~?
\item Qu'a en commun ce problème avec les bobines de câble du chapitre~9~?
\item Que changerait-il pour une plateforme installée dans un centre de données~?
\item En une phrase~: que garantit un \omterm{def:nw:private-connectivity-to-crypto-venues:endpoint}{point de terminaison privé}, et que ne garantit-il pas~?
\end{enumerate}
\end{problem}

\section{Questions d'entretien}

\begin{interviewq}[$\star$ \iqroles{développeur}]\label{iq:nw:private-connectivity-to-crypto-venues:1}
Quelle est la différence entre l'appairage de VPC et PrivateLink~?
\end{interviewq}

\begin{interviewq}[$\star\star$ \iqroles{développeur}]\label{iq:nw:private-connectivity-to-crypto-venues:2}
Comment vérifieriez-vous que votre connexion privée vers une plateforme ne franchit pas la frontière d'une \omterm{def:nw:trading-in-a-public-cloud:azid}{zone de disponibilité}~?
\end{interviewq}

\begin{interviewq}[$\star\star$ \iqroles{développeur, trader}]\label{iq:nw:private-connectivity-to-crypto-venues:3}
Une plateforme de crypto-actifs propose quatre paliers d'accès. Comment choisissez-vous, et que demandez-vous~?
\end{interviewq}

\begin{interviewq}[$\star\star$ \iqroles{développeur}]\label{iq:nw:private-connectivity-to-crypto-venues:4}
Qu'est-ce qui ajoute de la latence dans un chemin par \omterm{def:nw:private-connectivity-to-crypto-venues:endpoint}{point de terminaison privé}, et comment mesureriez-vous chaque élément~?
\end{interviewq}

\begin{interviewq}[$\star\star$ \iqroles{trader}]\label{iq:nw:private-connectivity-to-crypto-venues:5}
Un \omterm{def:nw:trading-in-a-public-cloud:pg}{groupe de placement} partagé avec la plateforme est-il équitable envers les autres intervenants du marché~?
\end{interviewq}

\begin{interviewq}[$\star\star\star$ \iqroles{développeur}]\label{iq:nw:private-connectivity-to-crypto-venues:6}
Concevez la connectivité vers une plateforme qui fonctionne dans deux zones d'une même région, à la fois pour la vitesse et pour la résilience.
\end{interviewq}
